% Movers for Message Passing
%
% Build:  latexmk -pdf movers.tex      (or: pdflatex movers.tex)
%
% Review comments: write  \RM{your note here}  anywhere in the text.
% They appear inline, in colour, bracketed.  Set \showrmfalse below to hide
% them all without deleting them.
%
\documentclass[11pt]{article}

\usepackage[T1]{fontenc}
\usepackage[utf8]{inputenc}
\usepackage{libertinus}
\usepackage[scaled=0.86]{beramono}
\usepackage{microtype}
\usepackage[a4paper,margin=1.15in]{geometry}
\usepackage{xcolor}
\usepackage{listings}
\usepackage{tcolorbox}
\usepackage{enumitem}
\usepackage{array}
\usepackage{booktabs}
\usepackage[hidelinks]{hyperref}
\usepackage{parskip}

\tcbuselibrary{skins,breakable}

% ---------------------------------------------------------------- colours ---
\definecolor{accent}{HTML}{1B4D8F}   % structure
\definecolor{leftmv}{HTML}{0F6E63}   % left movers
\definecolor{rightmv}{HTML}{9A3A12}  % right movers
\definecolor{neutralmv}{HTML}{6B5B8A}% non-movers
\definecolor{rulegrey}{HTML}{C9CFD8}
\definecolor{sunken}{HTML}{F2F4F7}
\definecolor{comment}{HTML}{6A737D}
\definecolor{rmcolor}{HTML}{B3261E}

% ------------------------------------------------------------ RM comments ---
\newif\ifshowrm \showrmtrue   % <-- \showrmfalse to hide every \RM{...}
\newcommand{\RM}[1]{%
  \ifshowrm{\color{rmcolor}\sffamily\small\,[RM: #1]\,}\fi}

  \newcommand{\CLAUDE}[1]{%
    \ifshowrm{\color{rmcolor}\sffamily\small\,[CLAUDE: #1]\,}\fi}
% ------------------------------------------------------------------ code ----
\newcommand{\code}[1]{\texttt{#1}}

\lstdefinelanguage{Verus}{
  morekeywords={fn,spec,proof,exec,open,closed,uninterp,pub,struct,enum,impl,trait,
    type,let,mut,ghost,tracked,requires,ensures,invariant,decreases,assert,assume,
    forall,exists,choose,match,if,else,while,loop,return,as,int,nat,bool,true,false,
    Self,self,use,mod,where,is,in,broadcast},
  sensitive=true,
  morecomment=[l]{//},
  morecomment=[s]{/*}{*/},
  morestring=[b]",
}
\lstset{
  language=Verus,
  basicstyle=\ttfamily\footnotesize,
  keywordstyle=\color{accent},
  commentstyle=\color{comment}\itshape,
  stringstyle=\color{leftmv},
  columns=fullflexible,
  keepspaces=true,
  showstringspaces=false,
  breaklines=true,
  breakatwhitespace=true,
  frame=single,
  rulecolor=\color{rulegrey},
  framesep=6pt,
  backgroundcolor=\color{sunken},
  xleftmargin=0pt,
  aboveskip=1.1em,
  belowskip=0.6em,
  captionpos=b,
  abovecaptionskip=4pt,
}
% caption above the listing, in the style of a figure label
\newcommand{\srcname}[1]{%
  \vspace{0.9em}\noindent{\sffamily\scriptsize\color{comment}\MakeUppercase{#1}}\par
  \vspace{-0.7em}}

% ------------------------------------------------------------ proof rules ---
\newtcolorbox{prooframe}[1]{
  enhanced, breakable,
  colback=white, colframe=rulegrey,
  boxrule=0.4pt, leftrule=2pt,
  borderline west={2pt}{0pt}{accent},
  arc=1pt, left=8pt, right=8pt, top=6pt, bottom=6pt,
  title={}, coltitle=black,
  before upper={\noindent{\ttfamily\scriptsize\color{accent}\textbf{#1}}\par\vspace{3pt}},
}
\newcommand{\reify}[1]{%
  \par\vspace{4pt}\noindent\textcolor{rulegrey}{\rule{\linewidth}{0.4pt}}\par
  \vspace{2pt}{\small\color{comment}#1}}

% ----------------------------------------------------------- trace strips ---
\newcommand{\chipR}[1]{\fcolorbox{rightmv}{white}{\color{rightmv}\ttfamily\scriptsize\vphantom{Ay}#1}}
\newcommand{\chipL}[1]{\fcolorbox{leftmv}{white}{\color{leftmv}\ttfamily\scriptsize\vphantom{Ay}#1}}
\newcommand{\chipN}[1]{\fcolorbox{neutralmv}{white}{\color{neutralmv}\ttfamily\scriptsize\vphantom{Ay}#1}}
\newenvironment{trace}
  {\par\vspace{0.7em}\noindent\hspace*{\fill}}
  {\hspace*{\fill}\par\vspace{0.7em}}

% ------------------------------------------------------------- structure ----
\usepackage{titlesec}
\titleformat{\section}{\sffamily\large\bfseries\color{black}}{\color{accent}\thesection}{0.7em}{}
\titleformat{\subsection}{\sffamily\normalsize\bfseries}{\thesubsection}{0.6em}{}
\titlespacing*{\section}{0pt}{1.8em}{0.6em}

\newcommand{\eyebrow}[1]{{\sffamily\scriptsize\color{accent}\MakeUppercase{#1}}\par\vspace{0.4em}}

\title{\vspace{-2.2em}\sffamily\bfseries Verifying Distributed Message Passing Protocols in Rust\vspace{-0.5em}}
\author{}
\date{}

\begin{document}
\maketitle
\vspace{-3.5em}

\noindent{\large How to prove a message-passing protocol by reasoning about one thread at a time.}

\vspace{1.2em}
\noindent\textcolor{rulegrey}{\rule{\linewidth}{0.4pt}}


% ===========================================================================
\section{The Problem}\label{sec:2pc}

We consider programs written in a small message-passing language. A program is a
fixed set of threads that share no memory and communicate only by passing
messages over \emph{channels}. A channel is created dynamically and has two
endpoints: a \emph{sender}, which may be held by several threads at once, and a
\emph{receiver}, which is held by exactly one. Sending is non-blocking; a thread
never waits to hand a message off. Receiving blocks until a message is
available.

Which message a receive returns is fixed by the channel's \emph{discipline}. We
consider two. Under \textsc{fifo}, messages are delivered in the order they were
sent, and we identify a channel with an ordered pair
\code{(source, destination)}, so a \textsc{fifo} channel has a single sender.
Under \textsc{bag}, a channel is an unordered pool and a receive returns any
message that has not yet been consumed; many threads may send on it.

We assume a Rust library that provides exactly this, and that the library itself
has been verified, so that its specification may be taken as given. On top of
it we also assume a verified \code{rpc} primitive, since request/response is the
dominant pattern in practice and we would like protocol code to use it directly.
Section~\ref{sec:rpc} shows how \code{rpc} is built and verified from the
message-passing operations; Section~\ref{sec:assumed} lists precisely what
remains assumed.

The question this note answers is: given such a library, how do we verify a
\emph{protocol} written on top of it?

As a running example, consider the two-phase commit protocol written on top of our message-passing
library. A
coordinator asks each participant whether it can commit; if every answer is yes,
it commits.

\srcname{Two-phase commit in Rust, abridged}
\begin{lstlisting}
// Coordinator
fn coordinator(peers: &[Peer]) -> bool {
    let mut all_yes = true;
    for p in peers {
        p.req.send(Prepare);              // nonblocking
        match p.rsp.recv() {              // blocks
            Vote(true)  => {}
            _           => all_yes = false,
        }
    }
    for p in peers { p.req.send(Decide(all_yes)); }
    all_yes
}

// Participant
fn participant(rx: Receiver<Msg>, tx: Sender<Msg>, vote: bool) {
    loop {
        match rx.recv() {                 // blocks
            Prepare   => tx.send(Vote(vote)),   // nonblocking
            Decide(b) => commit_or_abort(b),
            _         => {}
        }
    }
}
\end{lstlisting}

Both phases of the protocol reach a participant on the same channel, so a single
receive serves both: the participant does not know, before it looks, whether the
message in front of it is a request to vote or the announcement of the outcome.
That is worth noticing because it fixes how many interference points the
participant has. It has one per loop iteration --- the receive --- and not one
per protocol phase. The verified version in \code{src/twophase.rs} has exactly
this shape.

That receive does carry a proof obligation --- a \emph{yield invariant}, in the
sense of Section~\ref{sec:yield}, which the participant may assume when it wakes
up and must preserve at each of its own steps. For two-phase commit at this
level of abstraction the obligation is discharged trivially, because the
property we are after is a property of what is on the reply channels, and the
participant only ever appends to them. Section~\ref{sec:fanout} shows a version
of the same protocol whose yield invariant is not trivial, and
Section~\ref{sec:cr} one where establishing it is most of the work.

The property we want is:

\begin{quote}
If the coordinator commits, then every participant voted yes.
\end{quote}

Proving it directly in a verifier like Verus is quite complicated.
The coordinator and every participant run concurrently, so the proof must characterize every
state reachable under every interleaving of messages. An inductive invariant has to
describe, for an arbitrary point in the execution, which requests are in flight,
which have been consumed, which replies exist but have not yet been read, and
how those sets relate to the coordinator's loop counter. None of that appears in
the programmer's own reasoning, which is closer to:

\begin{quote}
The coordinator asks each participant in turn --- the order does not matter ---
and each participant answers with its own vote. When every answer is in, the
coordinator decides.
\end{quote}

The informal argument works for the following (informal) reasons.
First, the coordinator treats
\code{send} followed by \code{recv} as though it were an ordinary function call:
nothing observable happens in between. Second, the order in which unrelated
participants are asked does not matter.
Thus, correctness is based on sequential reasoning on the coordinator, with atomic
and independent calls to each participant.

The methodology described here is a way
to make both of those claims into checked proof obligations, so that the
resulting proof looks like the informal argument rather than like a case
analysis over interleavings.

% ===========================================================================
\section{Overview: Compositional verification}\label{sec:comp}

Before introducing our techniques, it is worth outlining how we verify concurrent 
software and what the difficulties are. 
We focus on safety properties, which are properties that must hold in every reachable state of the program.
The basic technique to prove safety properties is to find an inductive invariant, 
which is a property that holds in the initial state and is preserved by every step of the program.
The issue with concurrent programs is that the program's state space is huge, 
because it includes all possible interleavings of the threads.
A ``global'' invariant that describes the entire program state is often too complex to find or reason about:
it has to maintain the global state of all threads and the shared resources and this is usually complex.

Instead, we look for \emph{compositional} approaches.
A compositional approach allows us to reason about individual threads in isolation, 
and then combine those proofs to conclude properties about the whole program.


\subsection*{Owicki--Gries Reasoning}

The classical way to verify a concurrent program compositionally is to annotate
each thread with an assertion at every program point, as if it ran alone, and
add additional checks that ensure that the assertions remain valid under the
behavior of every other concurrently running thread. 
Then, we verify two different proof obligations:

\begin{itemize}[leftmargin=1.2em,itemsep=2pt,topsep=3pt]
\item \textbf{Sequential correctness.} Each thread's annotation is a valid
  Hoare-style proof of that thread in isolation.
\item \textbf{Non-interference.} Every assertion in every thread remains true under every
  atomic step of every other thread.
\end{itemize}

The second family typically leads to many checks. It has one obligation per
(assertion, foreign atomic step) pair, so it grows with the product of the
program sizes, and the assertions have to be phrased finely enough to be
preserved by whatever another thread might do next. 
For a protocol, this means reasoning about the network at the granularity of individual sends and receives.

One can simplify the form of the assertions by using a \emph{yield invariant}, 
which is a single predicate that every thread may assume at an interference point 
and must preserve at each of its own steps.

\subsection*{Ingredient 1: yield invariants}

A \emph{yield invariant} is a predicate over the shared state. Each
thread may assume it as a safe approximation of all concurrent activity in the program
on arrival at an interference point and must establish
that its own actions preserve the yield invariant.
The obligation is:

\begin{prooframe}{NON-INTERFERENCE}
Every atomic action of the program preserves the yield invariant.
\reify{One obligation per action of the protocol, rather than one per
(assertion, foreign step) pair.}
\end{prooframe}

The price is that the invariant must be strong enough to carry everything a
thread needs across an interference point, and weak enough that every action
preserves it. Choosing it well is the main creative act in these proofs; much of
Section~\ref{sec:net} is about picking a representation of the network that
makes such invariants easy to state and cheap to preserve.

\medskip

In principle, we should put yield invariants between any two atomic operations of each thread
since other threads may interleave at those points and change the shared state.
But we can often reduce the number of 
non-interference checks using a combination of two additional ingredients: 
reduction and refinement layers.

The idea of reduction is that some interleavings are redundant: 
they can be rearranged into an equivalent execution that ``moves'' a sequence of actions
of one thread next to each other and treats them as a single atomic action.
In this case, we do not need to put a yield invariant between those actions, 
since they are treated as a single atomic action without losing any generality.





\subsection*{Ingredient 2: reduction, via movers}

Lipton's insight is that the organization of a sequence of actions into a coarser
atomic block can be decided per action. An action is a
\emph{right mover} if, whenever it is immediately followed by an action of
another thread, the two can be swapped without changing the outcome; and a
\emph{left mover} if it can be swapped with an immediately preceding action of
another thread. 

In our setting, each thread works on its own local state and communicates by sending or receiving
messages.
Clearly, local computation in one thread 
does not interfere with other threads, 
so we can treat a sequence of local computation as a single atomic action.
Moreover, sending a message does not interfere with other threads; in particular,
sending enables other threads to receive that message, but it does not disable any other action.
Conversely, receiving a message does not interfere with other threads; it disables some actions, 
but it does not enable any other action.
In Lipton's terminology, sending is a left mover and receiving is a right mover.
Thus, a handler that receives a message, performs some local computation, and sends a reply has the shape
\chipR{recv}\,$\cdot$\,\chipN{compute}$\cdot$\,\chipL{send}
and can be treated as a single atomic action.

The reduction rule states the following:

\begin{prooframe}{REDUCE}
If a thread executes a sequence of actions consisting of zero or more right
movers, then at most one arbitrary action, then zero or more left movers, that
whole sequence may be treated as a \emph{single atomic action}.
\reify{Nothing in the executable code changes. What changes is where the proof
must consider interference: only between such blocks, not inside them.}
\end{prooframe}

Concretely, a thread that runs \chipR{recv}\,$\cdot$\,\chipN{compute}\,%
$\cdot$\,\chipL{send} could be intercepted before or after each action.
But reduction allows us to treat the whole sequence as a single atomic action, 
so we only have to consider interference before the receive and after the send.
This corresponds to our initial intuition that each participant is an atomic function call 
from the coordinator's point of view and we do not have to interleave between the receive and the send.



\subsection*{Ingredient 3: refinement layers}

The third ingredient attacks detail rather than interleaving. A \emph{layer} is
a complete program with its own set of atomic actions. A concrete layer
\emph{refines} an abstract one when every step of the concrete program is
matched by a step of the abstract program:

\begin{prooframe}{REFINE}
If the concrete layer refines the abstract layer, then any safety property
proved of the abstract layer holds of the concrete one.
\reify{The abstract layer's actions are typically coarser and nondeterministic:
they say \emph{that} something happened without saying exactly what. Callers
reason with the abstraction; the refinement obligation is discharged once.}
\end{prooframe}

Layers compose, so a realistic implementation can be connected to an obviously
correct specification through a chain of small steps rather than one large one.

\subsection*{The mechanism: ownership and tokenized state machines}

The three ingredients above are classical. What decides how expensive they are
is a representation choice: how the shared state is written down, and who is
entitled to change it.

The choice made here is to divide the state rather than share it. A
\emph{tokenized state machine} is a declaration of a transition system whose
state is split into named pieces. For each piece the Verus library generates a
\emph{token} type: a proof-mode value subject to Rust's ownership rules. A
thread may read or change a piece only by presenting the corresponding token,
and the declaration's invariants are proved inductive over its transitions once
and for all.

Two kinds of piece matter here.

\begin{itemize}[leftmargin=1.2em,itemsep=3pt,topsep=3pt]
\item An \textbf{exclusive} piece has one token, which cannot be duplicated. A
  thread holding it knows that no other thread can change that piece --- not
  because an invariant says so, but because the permission to change it is the
  token, and the thread has it.
\item A \textbf{persistent} piece has duplicable tokens that are never
  retracted. A fact recorded there is public and permanent, so it can be handed
  to another thread and will not become stale.
\end{itemize}

This is what makes the three ingredients cheap in what follows.
Non-interference for an exclusively owned piece needs no argument, because
interference on it is impossible. The mover conditions of
Section~\ref{sec:movers} become properties of the \emph{shape} of an operation
rather than claims about a particular protocol. And the invariant a thread must
carry across an interference point shrinks to what is genuinely shared, which
here is the record of what has been sent.

The price is that the state has to be divisible in the first place: the design
of Section~\ref{sec:net} is mostly the work of choosing a division under which
the facts a protocol needs are facts about pieces someone owns.

The code from here on is Rust with Verus annotations.
Appendix~\ref{sec:verus} collects the notation, including the syntax of a
tokenized state machine and a small self-contained example; a reader who knows
Rust but not Verus may want to read it first.

\subsection*{How they fit together here}

Reduction makes the atomic blocks in a protocol coarse enough that a thread's
code between two interference points looks like straight-line sequential code.
Yield invariants discharge what is left. Refinement layers let a component ---
in our case, a whole remote procedure call --- be replaced by a single action,
so its caller never sees the messages at all. Applied to two-phase commit, the
result is that the coordinator is verified as an ordinary sequential program
with an ordinary loop invariant, which is where we will arrive in
Section~\ref{sec:rpc}.

% ===========================================================================
\section{Modeling the network}\label{sec:net}

As a program executes, messages are sent and received. The ``shared state'' of a
message-passing program is the state of the network, and a proof of correctness
has to reason about it. How should it be represented?

The first question is not what to store but \emph{who owns it}. A single global
object holding the whole network is the obvious choice and the expensive one:
every invariant mentioning it must be re-established after every step of every
thread, and every fact a thread wants to carry across an interference point must
be re-derived from an explicit account of what other threads might have done.
That account is exactly the non-interference obligation of
Section~\ref{sec:comp}, and it reappears at every program point.

Instead we divide the network into pieces and give each piece an owner. The
representation has three parts, and each is chosen so that a fact about it is
stable for a structural reason rather than a proved one.

\srcname{src/tok.rs --- the three parts of the network state}
\begin{lstlisting}
#[sharding(map)]            pub sent:  Map<ChanId, Seq<M>>,
#[sharding(map)]            pub recvd: Map<ChanId, Seq<M>>,
#[sharding(persistent_set)] pub was_sent: Set<(ChanId, nat, M)>,
\end{lstlisting}

\textbf{The send history of a channel is owned by its sender.} Under per-link
FIFO a channel is identified by a source and a destination, so it has exactly
one sender. \code{sharding(map)} gives one token per channel, and a token cannot
be duplicated. A thread holding channel $c$'s \code{sent} token therefore knows
that nobody else appends to $c$ --- not because an invariant says so, but
because the permission to append is the token, and the thread has it.

\textbf{The consumption record of a channel is owned by its receiver.} Same
reasoning, resting on \code{Receiver} not being \code{Clone}.

\textbf{What was sent is public and permanent.} \code{was\_sent} records
``index $i$ of channel $c$ carries message $m$''. It is \emph{persistent}: its
tokens are duplicable and are never retracted. This is how a receiver learns
anything at all about a history it does not own. It cannot read the sender's
\code{sent} token, and it does not need to: a receive hands back a witness, and
a witness never goes stale.

The consequence is worth stating plainly, because it removes something that
would otherwise appear on every page. \emph{There is no environment step.} A
thread does not describe what other threads may have done and re-derive its
facts from that description. Facts about channels it owns cannot be disturbed;
facts about other channels are persistent and cannot be retracted; and the
protocol's invariant is global and inductive. Nothing has to be re-established
after an interference point, which is why the code that follows an interference
point in Section~\ref{sec:yield} is empty.

\subsection*{Naming channels}

Protocols need to know that distinct participants have distinct channels, and
that a request channel is never a reply channel. If a channel identifier is an
opaque number these are assumptions. We make them theorems instead by giving
identifiers structure.

\srcname{src/tok.rs --- channel names}
\begin{lstlisting}
pub struct ChanId { pub fam: nat, pub ix: Seq<int> }

pub open spec fn chan(fam: nat, ix: Seq<int>) -> ChanId { ChanId { fam, ix } }
\end{lstlisting}

A name is a family tag together with indices within that family: requests are
family $0$, replies family $1$, and the indices say which participant and which
call. Two names are equal exactly when their tags and indices are, which Verus
knows without being told. A protocol writes \code{req\_chan(j) = chan(0, seq![j])}
and \code{rsp\_chan(j,k) = chan(1, seq![j, k])}, and the distinctness facts it
needs follow by two small injectivity lemmas proved once in the library.

% ===========================================================================
\section{Movers, and why they need no annotation}\label{sec:movers}

This section comes before the first protocol because of what it does
\emph{not} produce: a protocol author writes no mover annotations, declares no
footprints, and proves no commutativity obligations. That is worth explaining
before the reader meets a protocol declaration and wonders what is missing from
it.

Reduction rests on Lipton's mover types. An action is a \emph{right mover} if it
commutes with any action of another thread that follows it, and a \emph{left
mover} if it commutes with any action of another thread that precedes it. A
sequence of actions executed by one thread that has the shape

\begin{trace}
\chipR{right movers}\;$\cdot$\;\chipN{one non-mover}\;$\cdot$\;\chipL{left movers}
\end{trace}

can be treated as a single atomic action: any interleaved execution can be
permuted into one in which the whole sequence runs uninterrupted.

For channels the assignment is forced, and it is the opposite of what one might
first guess.

\begin{prooframe}{SEND-L}
\textbf{\code{send} is a left mover.} Appending to a channel only ever
\emph{enables} other actions, never disables one, so \code{x; send} can always be
rewritten \code{send; x}. The converse fails: \code{send(m); recv(m)} cannot
become \code{recv(m); send(m)}, because the message is not there yet.
\end{prooframe}

\begin{prooframe}{RECV-R}
\textbf{\code{recv} is a right mover, and it blocks.} A blocking operation
violates the non-blocking side condition on left movers and can never be a left
mover. It is therefore the operation at which a thread yields.
\end{prooframe}

The consequence is the shape of a reduced atomic block:

\begin{trace}
\chipR{recv}\;$\cdot$\;\chipN{local compute}\;$\cdot$\;\chipL{send}\;$\cdot$\;\chipL{send}
\end{trace}

which is exactly the shape of an event handler. A handler loop needs one
interference point per iteration, at the receive, rather than one per channel
operation. Every protocol in Sections~\ref{sec:pingpong} to~\ref{sec:cr} has this
shape, and none of them says so anywhere in its text.

\subsection*{Why the mover types are automatic}

In a design where the network is one shared object, \textsc{send-l} is a claim
about a particular protocol, and it has to be proved for each one: that an
action reads and writes only within a declared footprint, that its gate is
preserved by unrelated appends, and that any two actions with disjoint
footprints commute. That is an obligation per action and, for commutativity, per
pair of actions.

Under the division of state described in
Section~\ref{sec:comp} the situation changes, because the two things a send does
are both invisible to other threads for reasons that have nothing to do with the
protocol. A send does exactly two things:

\begin{itemize}[leftmargin=1.2em,itemsep=2pt,topsep=3pt]
\item it changes a token that no other thread can hold, so no other thread's
  step can observe the change or be affected by it; and
\item it adds an element to a persistent set. Set union commutes with everything,
  and a persistent element is never removed, so this can only enable another
  action, never disable one.
\end{itemize}

Both bullets are exactly the two halves of \textsc{send-l}: appending cannot
disable another action, and it cannot be observed by one. Neither mentions the
protocol, the message, or the channel. The mover type is therefore fixed by the
\emph{shape} of the operation --- which token it consumes and which field it
extends --- and is settled once in the library rather than declared per protocol.
The same holds for \textsc{recv-r}, whose right-moverness comes from the receive
token being exclusively owned.

That is what \emph{automatic} means here, and it is worth being precise about
its limits: the reasoning above is a proof about the library's operations, not a
theorem Verus checks about all executions. What Verus checks is that a protocol
uses only those operations, and it enforces this by construction, since a thread
cannot touch a channel whose token it does not hold. What remains assumed is
Lipton's theorem, below.

This also gives a criterion for gates that can be checked by reading one
declaration. A gate preserves left-moverness when it mentions only protocol
constants, persistent facts, or state the thread exclusively owns. A gate that
mentions another thread's mutable state does not --- an action gated on
\emph{channel $c$ is empty} would be disabled by an unrelated append, and then
\code{send} would not be a left mover at all. Such a gate cannot be written
here: a gate is given the channel it is about and that channel's own history,
and has no way to name another. What an earlier design had to detect, this one
cannot express.

\subsection*{What is still assumed}

Lipton's theorem itself. That pairwise commutativity licenses the reordering is
a statement about all executions of the program, and there is no formal
execution semantics inside the Rust program to state it against. The two bullets
above are its side conditions, and they are discharged; the theorem that turns
them into atomicity is not, and Section~\ref{sec:assumed} says so.

% ===========================================================================
\section{A first protocol}\label{sec:pingpong}

The smallest protocol that uses every part of the method. Two parties and two
channels: \emph{A} sends \code{Ping} on the request channel, \emph{B} answers
with \code{Pong(v)} on the answer channel, and \emph{A} receives the answer and
knows the payload is well formed. The property we want is that last clause, and
the thing to watch for is that \emph{A}'s proof of it never mentions \emph{B}.
The code is \code{src/examples/pingpong.rs}.

\srcname{src/examples/pingpong.rs --- messages and channels}
\begin{lstlisting}
pub uninterp spec fn ok(v: u64) -> bool;      // what a good payload is

pub open spec fn ping_chan() -> ChanId { chan(0, seq![]) }
pub open spec fn pong_chan() -> ChanId { chan(1, seq![]) }

pub enum Msg { Ping, Pong(u64) }
\end{lstlisting}

The channels are built with \code{chan}, so \code{ping\_chan() != pong\_chan()}
follows from how they are named rather than being assumed;
\code{lemma\_chan\_distinct} discharges it.

\subsection*{The property, stated once}

Everything the protocol guarantees is one predicate on a message and the channel
it was sent on.

\begin{lstlisting}
pub open spec fn pong_ok(c: ChanId, m: Msg) -> bool {
    c == pong_chan() ==> (m is Pong && ok(m->Pong_0))
}
\end{lstlisting}

An answer is a well-formed \code{Pong}; nothing is claimed about the request
channel. This predicate will appear twice below, in two roles, and noticing that
they are the same predicate is most of understanding the method.

\subsection*{The protocol}

A protocol is an implementation of one trait. There is no state machine to
write: Section~\ref{sec:machine} describes the one every protocol shares.

\srcname{src/examples/pingpong.rs --- the whole specification}
\begin{lstlisting}
pub struct PingPong;

impl NetInv<Msg> for PingPong {
    open spec fn gate(c: ChanId, s: Seq<Msg>, m: Msg) -> bool { pong_ok(c, m) }

    open spec fn wit_inv(c: ChanId, m: Msg) -> bool { pong_ok(c, m) }

    open spec fn deliverable_at(v: Seq<Msg>, i: nat) -> bool { fifo_deliverable(v, i) }

    open spec fn extra(sent: Map<ChanId, Seq<Msg>>) -> bool { true }

    proof fn lemma_gate_gives_inv(c: ChanId, s: Seq<Msg>, m: Msg) { }
    proof fn lemma_extra_init(chans: Set<ChanId>) { }
    proof fn lemma_extra_alloc(sent: Map<ChanId, Seq<Msg>>, c: ChanId) { }
    proof fn lemma_extra_preserved(...) { }
}
\end{lstlisting}

Four definitions and four proofs, and the proofs are empty. Taking them in turn:

\textbf{\code{gate}} is the condition under which a send must not \emph{fail}.
This is Civl's $\rho$, and it is distinct from blocking: a send whose gate does
not hold is a bug in the program, whereas an operation with no available step
merely waits. Here, putting a malformed answer on the answer channel is a bug.
The obligation appears at every call site of \code{send}.

\textbf{\code{wit\_inv}} is what may be concluded about a message that
\emph{was} sent. This is the yield invariant: a thread arriving at an
interference point may assume it of anything on any channel. Note that it is
\code{pong\_ok} again --- the gate says what may be placed on a channel, and the
invariant says everything there satisfies it. They are one predicate read at two
moments, which is the usual case and is why
\code{lemma\_gate\_gives\_inv} is empty. Section~\ref{sec:machine} says what
happens when they differ.

\textbf{\code{deliverable\_at}} is the delivery discipline: given what has been
consumed on a channel, which index may arrive next. Per-link FIFO pins it to the
current length. Section~\ref{sec:unreliable} gives the other case.

\textbf{\code{extra}} is for a guarantee no statement about a single message can
express --- an ordering \emph{between} messages, say. Most protocols leave it
\code{true}; Section~\ref{sec:machine} returns to the one that does not.

\subsection*{The two parties}

\srcname{src/examples/pingpong.rs}
\begin{lstlisting}
pub fn responder(rx: &Receiver<Msg>, tx: &Sender<Msg>, v: u64, ...)
    requires ok(v), rx.id@ == ping_chan(), tx.id@ == pong_chan(), ...
{
    let _m = recv::<Msg, PingPong>(rx, Tracked(inst), Tracked(rtok));
    let _w = send::<Msg, PingPong>(tx, Msg::Pong(v), Tracked(inst), Tracked(stok));
}

pub fn initiator(tx: &Sender<Msg>, rx: &Receiver<Msg>, ...) -> (answer: u64)
    ensures ok(answer),
{
    let _w = send::<Msg, PingPong>(tx, Msg::Ping, Tracked(inst), Tracked(stok));
    let m = recv_learn::<Msg, PingPong>(rx, Tracked(inst), Tracked(rtok));
    match m {
        Msg::Pong(v) => v,
        Msg::Ping    => { proof { assert(false); } 0 }
    }
}
\end{lstlisting}

Both are ordinary sequential Rust. Three things are worth pointing at.

\textbf{The responder has one interference point.} Its body is a blocking
receive followed by a send --- \chipR{R}$\,\cdot\,$\chipL{L} in the terms of
Section~\ref{sec:movers} --- so the two operations form one atomic block and the
only place another thread runs is at the receive. Nothing in the code says so.

\textbf{The gate obligation appears at the send.} \code{send} requires
\code{PingPong::gate} of the channel and message, which here reduces to
\code{ok(v)}, discharged by the responder's precondition. Weaken the gate and
the protocol's guarantee goes with it; the development includes that experiment.

\textbf{The initiator's proof does not mention the responder.} It sends, yields,
receives, and concludes \code{ok(answer)}. The step that gets it there is
\code{recv\_learn}, which receives a message and converts the accompanying
witness into \code{wit\_inv}. Replacing it with a plain \code{recv} makes the
postcondition fail. \emph{A} reasons only about what the invariant says of
anything on that channel.

The \code{Ping} branch proves \code{false}: the invariant says the answer
channel carries only \code{Pong}s, so receiving a \code{Ping} there cannot
happen. That branch exists because Rust wants the match to be exhaustive, not
because the protocol admits it.

Every function that touches the network carries two extra arguments --- the
state machine instance and the token for the channel it uses --- together with
preconditions relating them. They are mechanical, and they are the main
notational cost of the approach.

% ===========================================================================
\section{The network state machine}\label{sec:machine}

Section~\ref{sec:pingpong} gave the four definitions a protocol supplies and
said what each means. This section is the other half: the single state machine
that gives them their meaning. It is ordinary library code, in
\code{src/tok.rs}, and it is shared by every protocol.

\srcname{src/tok.rs --- the state}
\begin{lstlisting}
tokenized_state_machine!{
    NetSM<M, Inv: NetInv<M>> {
        fields {
            #[sharding(map)]            pub sent:  Map<ChanId, Seq<M>>,
            #[sharding(map)]            pub recvd: Map<ChanId, Seq<M>>,
            #[sharding(persistent_set)] pub was_sent: Set<(ChanId, nat, M)>,
            #[sharding(variable)]       pub next: nat,
            #[sharding(constant)]       pub inv: PhantomData<Inv>,
        }
\end{lstlisting}

The first three fields are the representation of Section~\ref{sec:net}:
\code{sent} and \code{recvd} divided per channel and exclusively owned,
\code{was\_sent} persistent. \code{next} is the allocator used by \code{mint}
(Section~\ref{sec:rpc}). The last field carries the protocol as a type
parameter and holds no information; it is the idiom \code{vstd}'s own
\code{rwlock} uses for the same purpose.

\subsection*{Where each part of a protocol is used}

\srcname{src/tok.rs --- the transitions}
\begin{lstlisting}
transition!{
    do_send(c: ChanId, s: Seq<M>, m: M) {
        remove sent -= [c => s];
        require(Inv::gate(c, s, m));            // the protocol's gate
        add    sent += [c => s.push(m)];
        add    was_sent (union)= set { (c, s.len(), m) };
    }
}

transition!{
    do_recv(c: ChanId, r: Seq<M>, i: nat, m: M) {
        remove recvd -= [c => r];
        have   was_sent >= set { (c, i, m) };
        require(Inv::deliverable_at(r, i));     // the protocol's discipline
        add    recvd += [c => r.push(m)];
    }
}
\end{lstlisting}

The \code{remove} in each says a send consumes the channel's send token and a
receive its consumption token, so only a channel's sender may send on it and
only its single consumer may receive. The \code{have} in \code{do\_recv}
requires a witness for the message being taken: you can only receive what was
sent, and this is where that is enforced. The two \code{require}s are the only
places a protocol's own definitions enter.

The invariants are what a thread may assume:

\srcname{src/tok.rs --- the invariants}
\begin{lstlisting}
#[invariant]                            // a witness describes the history
pub spec fn agree(&self) -> bool { ... }

#[invariant]                            // the protocol's guarantee
pub spec fn protocol_inv(&self) -> bool {
    forall|c, i, m| self.was_sent.contains((c, i, m)) ==> Inv::wit_inv(c, m)
}

#[invariant]                            // the protocol's extra guarantee
pub spec fn protocol_extra(&self) -> bool { Inv::extra(self.sent) }

#[invariant]                            // nothing at or above the counter exists
pub spec fn unallocated(&self) -> bool { ... }
\end{lstlisting}

\code{protocol\_inv} is stated over \code{was\_sent}, which only ever grows, so
no other thread can falsify it: another thread may add a witness but cannot
remove one, and adding one preserves a claim about all existing witnesses
provided the new one satisfies it too. That is exactly what the gate guarantees,
and it is the whole of the proof the machine demands of a send:

\begin{lstlisting}
#[inductive(do_send)]
fn do_send_inductive(pre: Self, post: Self, c: ChanId, s: Seq<M>, m: M) {
    Inv::lemma_gate_gives_inv(c, s, m);
    Inv::lemma_extra_preserved(pre.sent, c, s, m);
    ...
}
\end{lstlisting}

A protocol's obligations are therefore small and fixed. \code{gate} must be
strong enough to establish \code{wit\_inv} for the message it admits ---
immediate when they are the same predicate, which is the usual case. And
\code{extra} must hold initially and be preserved by a send and by the creation
of a channel.

Finally, how a thread uses the invariant:

\srcname{src/tok.rs --- reading the invariant off a witness}
\begin{lstlisting}
property!{
    learn(c: ChanId, i: nat, m: M) {
        have was_sent >= set { (c, i, m) };
        assert(Inv::wit_inv(c, m));
    }
}
\end{lstlisting}

A \code{property!} changes nothing. It says: hold a witness, and you may
conclude \code{wit\_inv} of what it names. The \code{assert} is proved from the
invariant, once, in the library. \code{recv\_learn} is a receive followed by
this.

\subsection*{When the gate and the guarantee differ}

They coincide whenever the gate constrains the message alone. They cannot
coincide when a protocol's guarantee is about \emph{pairs} of messages, because
a witness names one message. \code{heartbeat.rs} is the case: each beat's
sequence number must exceed every earlier one on the link. Its \code{wit\_inv}
is \code{true} and the guarantee lives in \code{extra}:

\srcname{src/examples/heartbeat.rs}
\begin{lstlisting}
open spec fn gate(c: ChanId, s: Seq<Beat>, m: Beat) -> bool {
    forall|x: int| 0 <= x < s.len() ==> (#[trigger] s[x]).seq < m.seq
}
open spec fn wit_inv(c: ChanId, m: Beat) -> bool { true }
open spec fn extra(sent: Map<ChanId, Seq<Beat>>) -> bool {
    sent.dom().contains(link()) ==>
        forall|x: int, y: int| 0 <= x < y < sent[link()].len()
            ==> (#[trigger] sent[link()][x]).seq < (#[trigger] sent[link()][y]).seq
}
\end{lstlisting}

Note that the gate reads the history \code{s}, which is what \code{wit\_inv}
cannot do. The three \code{extra} lemmas are then not empty, and they are one of
the two places in the example protocols where a protocol writes a real proof
about its own invariant; the other is the journal ordering in
\code{leaselock.rs}.

All three of these read one channel. A guarantee that relates a message on one
channel to a message on another is outside what any of them can state, and needs
the separate mechanism of Section~\ref{sec:provenance}.

\subsection*{What a protocol does not write}

No state machine. No mover annotations, footprints or commutativity obligations
--- Section~\ref{sec:movers} gave the reason. No description of what other
threads may do. And no ghost state of its own: \code{sent}, \code{recvd},
\code{was\_sent} and the allocator are all the network is.

A protocol that needs ghost state beyond that --- a phase counter, a record of
local decisions --- declares a \emph{separate} state machine for it. The two are
independent instances, so a token of one cannot affect the other and no
obligation relates them, exactly as for two protocols sharing a thread
(Section~\ref{sec:threads}). The case this does not cover is protocol state that
must be related to network state; that has to live in one machine, and is not
expressible.
% ===========================================================================
\section{Interference points}\label{sec:yield}

\srcname{src/tok.rs}
\begin{lstlisting}
pub fn interference_point() { }
\end{lstlisting}

It takes no ghost state, promises nothing, and is not trusted. Marking a yield
is documentation, not an obligation, for the reason given at the end of
Section~\ref{sec:net}: there is nothing to re-establish. The heartbeat example
makes the point concretely.

\srcname{src/examples/heartbeat.rs}
\begin{lstlisting}
pub fn beat_and_wait(tx: &Sender<Beat>, seq: u64, ...)
    requires
        forall|x: int| 0 <= x < old(tok).value().len()
            ==> old(tok).value()[x].seq < seq,
    ensures
        final(tok).value().last() == (Beat { seq }),
{
    let _w = send::<HbTok>(tx, Beat { seq }, Tracked(inst), Tracked(tok));
    interference_point();
    // Nothing to prove here. Nobody else could have appended.
}
\end{lstlisting}

The thread still holds the link's \code{sent} token after the yield, so what it
wrote is still the last thing there. The conclusion is available because the
permission was never given away.

\subsection*{Where a protocol's definitions are checked}\label{sec:coherence}

A protocol states its gate, its guarantee and its delivery discipline once each,
and the machine of Section~\ref{sec:machine} is the only thing that reads them.
There is no second copy to keep consistent: \code{Inv::gate} appears in exactly
one \code{require}, \code{Inv::deliverable\_at} in one other, and
\code{Inv::wit\_inv} in one invariant. The obligations relating them ---
principally that the gate establishes the guarantee --- are the trait's proof
members, discharged per protocol and checked by the machine's own inductive
proof.

% ===========================================================================
\section{Remote procedure calls}\label{sec:rpc}

A thread that makes a request and waits for the answer writes \code{send} then
\code{recv}. That is a left mover followed by a right mover --- the wrong shape.
The two cannot be one atomic block, so the caller yields in the middle of what
it thinks of as a function call, and every invariant it holds must survive that.

The fix is Civl's: strengthen the gate of the receive until the operation cannot
block. A non-blocking receive is a left mover, and the pair becomes
\code{L}$\cdot$\code{L}, one block with no interference point inside it.

\begin{prooframe}{RECV-ABS}
\textbf{A receive whose gate guarantees the message is already present is a left
mover.} It cannot block, so the non-blocking side condition is satisfied.
\reify{\code{recv\_abs} in \code{src/tok.rs}. Its executable behaviour is the
same blocking receive; only the specification is stronger.}
\end{prooframe}

The gate is a witness. The caller must hold a persistent token saying that the
reply was sent, at the index it is about to consume.

\srcname{src/tok.rs --- the abstracted receive}
\begin{lstlisting}
pub fn recv_abs<T: DetDelivery>(r: &Receiver<T::M>, ..., Tracked(w): Tracked<&T::Wit>)
    -> (m: T::M)
    requires
        w.element().0 == r.id@,
        T::deliverable_at(old(tok).value(), w.element().1),
    ensures
        m == w.element().2,
\end{lstlisting}

Two things follow that the earlier design had to arrange by other means.

\textbf{Delivery is deterministic without freshness.} The witness names an
index, and the agreement invariant of Section~\ref{sec:net} says an index
determines the message. An earlier version needed each call to mint a
single-use channel and a separate lemma about singleton delivery to know
\emph{which} message came back. Knowing the index is enough.

\textbf{The requirement is now visible in the type.} \code{recv\_abs} is
available only for \code{T: DetDelivery}, a protocol that can prove at most one
index is deliverable at a time. An unreliable link cannot: it may hand back any
index at any moment. So \code{rpc} is unavailable on such a link, which is
correct and was previously not enforced anywhere.

\subsection*{Executing the handler at the call site}

Where does the caller's witness come from? From performing the handler's action
itself. This is Civl's synchronised asynchronous call: a left-moving handler
whose request has already been sent may be executed at the call site.

\srcname{src/tok.rs}
\begin{lstlisting}
pub proof fn absorb_handler<M, Inv: NetInv<M>>(
    tracked inst: &NetSM::Instance<M, Inv>,
    tracked reply_cap: NetSM::sent<M, Inv>,
    reply: M,
) -> (tracked r: (NetSM::sent<M, Inv>, NetSM::was_sent<M, Inv>))
{
    let tracked (a, b) = inst.do_send(reply_cap.key(), reply_cap.value(), reply, reply_cap);
    (a.get(), b.get())
}
\end{lstlisting}

This is a \code{proof fn}, not an axiom. The caller genuinely performs the
machine's own send transition, using the reply channel's send token, which it
holds because it created the channel. The obligation that remains --- that the
caller and the real handler do not \emph{both} send the reply --- is discharged
by Rust: \code{rpc} takes the reply capability by value and consumes it, because
a reply channel is used once.

\srcname{src/tok.rs --- the call, verified}
\begin{lstlisting}
pub fn rpc<T: DetDelivery>(
    req_s: &Sender<T::M>, req_m: T::M,
    rsp_r: &Receiver<T::M>, reply: Ghost<T::M>,
    Tracked(inst): ..., Tracked(req_tok): ...,
    Tracked(reply_cap): Tracked<T::Sent>,      // consumed
    Tracked(rsp_tok): ...,
) -> (m: T::M)
    ensures m == reply@,
{
    let _w = send::<T>(req_s, req_m, Tracked(inst), Tracked(req_tok));
    let tracked w2;
    proof { let tracked (_s, wit) = absorb_handler::<T>(inst, reply_cap, reply@); w2 = wit; }
    recv_abs::<T>(rsp_r, Tracked(inst), Tracked(rsp_tok), Tracked(&w2))
}
\end{lstlisting}

\subsection*{Creating channels}

A caller that mints a reply channel per call needs the channel to be fresh.
Freshness used to be asserted by the allocator's specification, which nobody
could check. It is now proved.

\srcname{src/examples/twophase.rs --- the allocator's invariant}
\begin{lstlisting}
#[invariant]
pub spec fn unallocated(&self) -> bool {
    forall|j: int, k: nat| k >= self.next
        ==> !self.sent.dom().contains(rsp_chan(j, k)) && ...
}
\end{lstlisting}

The allocator is an exclusively owned counter, and the invariant says nothing at
or above it exists yet. Adding a key to a sharded map carries an obligation that
the key is absent, and that obligation is discharged from this invariant. The
library's \code{mint} is then a verified function; the only trusted part is a
routine that fabricates two endpoint handles for a channel identifier and
touches no ghost state.

\subsection*{Two-phase commit, written as a sequential program}

\srcname{src/examples/twophase.rs --- the coordinator}
\begin{lstlisting}
while i < reqs.len() && all_yes
    invariant all_yes ==> forall|j: int| 0 <= j < i ==> vote(j),
{
    let (reply_tx, reply_rx, Tracked(reply_cap), Tracked(mut reply_rcv)) =
        mint::<TpcTok>(i, Tracked(inst), Tracked(next_tok));
    let m = rpc::<TpcTok>(&reqs[i], Msg::Prepare(Ghost(reply_rx.id@)),
                          &reply_rx, Ghost(Msg::Vote(vote(i as int))), ...);
    match m { Msg::Vote(b) => if !b { all_yes = false; }, _ => all_yes = false, }
    i = i + 1;
}
\end{lstlisting}

The loop has no interference point: each iteration mints a channel and makes one
call which, being \code{L}$\cdot$\code{L}, is atomic. The correctness argument
is the loop invariant, and it is the argument a programmer would give.

% ===========================================================================
\section{Fan-out two-phase commit}\label{sec:fanout}

The coordinator above asks one participant at a time. A deployed one broadcasts
and then gathers, stopping early on the first refusal. There is no remote call
here, so nothing hides the interference: every collect is an interference point,
and other participants reply while the coordinator is blocked.

\srcname{src/examples/twophase\_fanout.rs --- broadcast, then gather}
\begin{lstlisting}
while i < reqs.len() { ... send::<TpcfTok>(&reqs[i], FMsg::Prepare, ...); ... }

while k < rsps.len() && all_yes
    invariant all_yes ==> forall|j: int| 0 <= j < k ==> vote(j),
{
    let (m, w) = recv::<TpcfTok>(&rsps[k], Tracked(inst), Tracked(&mut rtok));
    proof {
        TpcfTok::learn(inst, w.borrow());
        assert(is_rsp(cid, k as int));
    }
    match m { FMsg::Vote(b) => if !b { all_yes = false; }, _ => all_yes = false, }
    k = k + 1;
}
\end{lstlisting}

The broadcast loop contains no interference point at all: every operation in it
is a send, and sends are left movers, so the whole fan-out is one block. The
gather loop has one per iteration.

What carries the proof is the yield invariant, declared as one line:

\begin{lstlisting}
wit_inv(c, m) { forall|j: int| is_rsp(c, j) ==> m == FMsg::Vote(vote(j)) }
\end{lstlisting}

Participant $j$'s reply channel carries nothing but participant $j$'s vote.
Where the remote-call version learned a participant's vote by executing that
participant's action, this one learns it from the invariant, through
\code{learn}. Both prove the same safety property by different routes, which is
a reasonable illustration of what the two techniques each buy.

\subsection*{What this version still does not do}

It gathers in a fixed order: participant $0$, then $1$, and so on. That is why a
plain loop invariant suffices. A coordinator that accepts the first $m$ of $n$
replies in whatever order they arrive needs an induction over the set of
outstanding replies, which is what inductive sequentialization provides and this
development does not have.

% ===========================================================================
\section{Chang--Roberts leader election}\label{sec:cr}

A ring of $n$ nodes with distinct identifiers. Node $i$ receives on link $i-1$
and sends on link $i$; on receiving a value $v$ it forwards if $v > \code{id}(i)$,
declares itself leader if $v = \code{id}(i)$, and otherwise drops the message. The
safety property is that only the node with the maximum identifier can declare
itself leader.

The interest here is that the invariant is genuinely about the protocol rather
than about the framework.

\srcname{src/examples/chang\_roberts.rs --- the ring invariant}
\begin{lstlisting}
pub open spec fn msg_ok(j: int, m: Elect) -> bool {
    &&& 0 <= m.origin@ < n_nodes()
    &&& m.val as int == id(m.origin@)
    &&& j == (m.origin@ + m.hops@) % n_nodes()
    &&& forall|k: int| 0 <= k <= m.hops@
            ==> m.val as int >= id((m.origin@ + k) % n_nodes())
}
\end{lstlisting}

A message on link $j$ carries its originator's identifier, has travelled
\code{hops} steps from that originator, and dominates every node it has passed
through. \code{origin} and \code{hops} are \code{Ghost}: they exist in the proof
and are erased from the running program.

A node's whole step is \code{R}$\cdot$\code{N}$\cdot$\code{L}: one blocking
receive, a comparison, and at most one send. One interference point. When the
value equals the node's own identifier, the invariant says the message has
travelled a full circle and dominated everything on the way, which is the safety
property. The supporting arithmetic --- that a full circle covers every node ---
is about modular arithmetic and identifiers, and is independent of the framework.

% ===========================================================================
\section{Facts that span channels}\label{sec:provenance}

Everything a protocol has said so far is confined to one channel. A gate is
given the history of the channel being written and the message being placed on
it; \code{wit\_inv} is given even less; \code{extra} ranges over the send
histories but is established one send at a time, by a thread that holds only the
channel it is writing. This is what makes the reasoning local, and it is also a
real limit. Some properties are about a message on one channel and a message on
another, and no amount of strengthening a gate will express them.

The lease lock is the standard example. It is written in
\code{src/examples/leaselock.rs} as three kinds of process exchanging
asynchronous messages: a lock server that hands out increasing lease tokens, one
per request; a storage node that accepts a write only if its token and sequence
number strictly exceed every write already accepted, appending accepted writes
to a journal; and concurrent writers that acquire a lease and then write. Lease
expiry is modelled by nondeterminism rather than by time. A writer may be
granted a lease at any moment, so an earlier writer's token can be superseded
while that writer still believes it holds the lock, and its next write is
refused. Nothing in the writer's code detects this, which is the situation a
fencing token exists to handle.

\srcname{src/examples/leaselock.rs --- channels}
\begin{lstlisting}
pub open spec fn acq_req(w: int) -> ChanId { chan(0, seq![w]) } // writer  -> server
pub open spec fn acq_rsp(w: int) -> ChanId { chan(1, seq![w]) } // server  -> writer
pub open spec fn wr_req(w: int)  -> ChanId { chan(2, seq![w]) } // writer  -> storage
pub open spec fn wr_rsp(w: int)  -> ChanId { chan(3, seq![w]) } // storage -> writer
pub open spec fn journal()       -> ChanId { chan(4, seq![]) }  // storage -> the record
\end{lstlisting}

The safety property proved is write serialization --- the journal is strictly
ordered by (token, sequence) --- and it is the top-level theorem of the
corresponding Leslie model. It is deliberately \emph{not} mutual exclusion: two
writers may simultaneously believe they hold the lease, and the protocol is
still correct. That part needs nothing new. The fencing check is a gate on
\code{journal()}, the storage node is the only sender there and so owns that
history, and the ordering is a fact about state it holds.

What does need something new is the claim that a write's token was issued by the
server at all. The token arrives on \code{wr\_req(w)}, the grant was sent on
\code{acq\_rsp(w)}, and a gate on the first cannot see the second. A forged
token and a real one look identical.

\subsection*{Provenance}

The mechanism is to make a send point at a message already sent elsewhere, and
to require the sender to hold the witness for it. Since witnesses are obtained
only by sending or receiving, a participant cannot point at a message it never
saw. Four definitions carry this, all optional and all defaulting to nothing:

\srcname{src/tok.rs --- the provenance interface}
\begin{lstlisting}
// Does sending m on c require pointing at a message sent somewhere else?
open spec fn needs_cause(c: ChanId, m: M) -> bool { false }

// Are the messages in `causes` acceptable justification for m on c?
open spec fn caused_by(c: ChanId, m: M, causes: Set<(ChanId, nat, M)>) -> bool { false }

// What a thread may conclude about m from the fact that a justification exists.
open spec fn cause_gives(c: ChanId, m: M) -> bool { true }

proof fn lemma_cause_gives(c: ChanId, m: M, causes: Set<(ChanId, nat, M)>)
    requires Self::needs_cause(c, m), Self::caused_by(c, m, causes),
             forall|e: (ChanId, nat, M)| causes.contains(e) ==> Self::wit_inv(e.0, e.2),
    ensures  Self::cause_gives(c, m);
\end{lstlisting}

A message with \code{needs\_cause} cannot go through \code{send}, whose
precondition denies it. It goes through \code{send\_caused}, which additionally
takes a witness and checks \code{caused\_by} against the message that witness
names, or through \code{send\_general}, which takes a whole set of them. There
is one send transition, and it carries the justification:

\srcname{src/tok.rs --- the transition}
\begin{lstlisting}
transition!{
    do_send(c: ChanId, s: Seq<M>, m: M, causes: Set<(ChanId, nat, M)>) {
        remove sent -= [c => s];
        have   was_sent >= (causes);
        require(Inv::gate(c, s, m));
        require(Inv::needs_cause(c, m) ==> Inv::caused_by(c, m, causes));
        add    sent += [c => s.push(m)];
        add    was_sent (union)= set { (c, s.len(), m) };
    }
}
\end{lstlisting}

The \code{have} is what does the work. It is a guard on a token the sender must
already own, so the justification cannot be invented. The corresponding
invariant says every message that needed one has one:

\srcname{src/tok.rs --- the invariant}
\begin{lstlisting}
#[invariant]
pub spec fn provenance(&self) -> bool {
    forall|c: ChanId, i: nat, m: M|
        (#[trigger] self.was_sent.contains((c, i, m))) && Inv::needs_cause(c, m)
            ==> exists|causes: Set<(ChanId, nat, M)>|
                    causes.subset_of(self.was_sent) && Inv::caused_by(c, m, causes)
}
\end{lstlisting}

This is inductive for the same reason the rest of the model is stable under
interference: \code{was\_sent} only grows. A justification found once cannot be
removed by a later step of any thread, so a message established as justified
stays justified.

In the lease lock the definitions read:

\srcname{src/examples/leaselock.rs}
\begin{lstlisting}
open spec fn needs_cause(c: ChanId, m: Msg) -> bool {
    c.fam == 2 && c.ix.len() == 1 && m is Write
}
open spec fn caused_by(c: ChanId, m: Msg, causes: Set<(ChanId, nat, Msg)>) -> bool {
    exists|j: nat| causes.contains((acq_rsp(c.ix[0]), j, Msg::Granted(m->Write_0)))
}
open spec fn cause_gives(c: ChanId, m: Msg) -> bool {
    Self::needs_cause(c, m) ==> m->Write_0 != 0
}
\end{lstlisting}

Note that \code{caused\_by} derives the expected channel from \code{c} rather
than quantifying over writers. A writer holding some other writer's grant cannot
use it, because the channel the justification must have come from is determined
by the channel the write is going out on.

\subsection*{Reading it back}

A thread that receives a write holds a witness for that message and nothing
else. It recovers the cross-channel fact through a property of the machine:

\srcname{src/tok.rs --- reading provenance off a witness}
\begin{lstlisting}
property!{
    learn_cause(c: ChanId, i: nat, m: M) {
        have was_sent >= set { (c, i, m) };
        require(Inv::needs_cause(c, m));
        birds_eye let ws = pre.was_sent;
        assert(Inv::cause_gives(c, m)) by {
            let causes = choose|causes: Set<(ChanId, nat, M)>|
                causes.subset_of(ws) && Inv::caused_by(c, m, causes);
            Inv::lemma_cause_gives(c, m, causes);   // (details elided)
        };
    }
}
\end{lstlisting}

Two things here deserve comment.

The \code{birds\_eye} binding reads the whole record, which no thread owns. That
is sound only because \code{was\_sent} is monotone, for exactly the reason given
above; the same binding on \code{sent} would be unsound, since another thread
can extend a history at any time.

The second is why \code{cause\_gives} exists at all, and it is the part that is
easy to get wrong. The justifying message is reached by an existential, and the
record it is quantified over is not something the caller can name. If the
property concluded \code{exists causes. causes \(\subseteq\) ws \&\& caused\_by(...)}
directly, then what the caller receives is that statement with the record
projected out, which for any concrete \code{caused\_by} is trivially satisfiable
--- one need only choose a \code{d} and \code{m2} of the right shape. The whole
chain would verify and establish nothing. \code{cause\_gives} is the part that
survives the projection: a predicate in \code{c} and \code{m} alone, whose proof
was allowed to use the justification. Its content comes from
\code{wit\_inv(d, m2)}, available because everything ever sent satisfies the
protocol's guarantee.

In the lease lock, the server's gate guarantees that an issued token is nonzero,
so \code{cause\_gives} says a write's token is nonzero --- a fact about the write
that the writer cannot manufacture. The storage node learns it on receipt, and
it feeds the gate on \code{journal()}, so the end-to-end statement is that every
write in the journal carries a token the server actually issued.

Four checks confirm the mechanism is not decorative. A writer that sends
\code{Write(999, ...)} after being granted a different token fails the
precondition of \code{send\_caused}; a writer that reaches for plain \code{send}
fails its \code{!needs\_cause} precondition; a storage node that omits the
\code{learn\_cause} call can no longer discharge the journal's gate; and
deleting the call from an otherwise correct proof makes it fail, which is what
rules out the vacuous version. The first is kept as
\code{counterexamples/forged\_lease.rs}.

\subsection*{What this does not reach}

Provenance relates a message to \emph{some} earlier message. It does not by
itself give an ordering between messages sent to different participants, so
global monotonicity of the server's grants --- that tokens handed to different
writers increase --- is still not expressible this way. Write serialization does
not need it, since it follows from the fencing gate alone, and the storage node
proves that without reference to the server or to any other writer.

% ===========================================================================
\section{Unreliable links and unordered delivery}\label{sec:unreliable}

Two variations that look similar and are not. It is worth separating them,
because conflating them costs either soundness or expressiveness.

\begin{itemize}[leftmargin=1.2em,itemsep=3pt,topsep=3pt]
\item \textbf{Who may send} determines whether a send history can be
  exclusively owned, and therefore whether \code{send} is a left mover.
\item \textbf{The delivery discipline} determines what a receiver may be handed.
  It is a property of the receive side alone.
\end{itemize}

\subsection*{An unreliable link}

Loss and duplication are of the second kind. A lossy link has one sender and one
receiver --- it is a datagram link, not a shared mailbox --- so ownership is
untouched and the entire difference from per-link FIFO is one line:

\begin{lstlisting}
deliverable(v, i) { true }        // any index, any number of times
\end{lstlisting}

Loss needs no modelling at all: nothing in this development ever forces a
delivery, so a message that never arrives leaves its receiver blocked.
Duplication is the dropped index constraint, and the persistent witness is what
makes re-delivery expressible --- a witness is duplicable and never retracted,
so handing the same one back twice is what a duplicating network does.

What survives unchanged is any invariant of the form \emph{everything ever sent
satisfies $P$}, together with the step from a message in a receiver's hand back
to the send history. Both are statements about \code{sent}, which duplication
and loss do not touch. \code{receive\_one} returns \code{ok(p.v)} on an unreliable
link with no additional work.

What does not survive is counting, and the development demonstrates the
difference rather than asserting it. The same claim --- that two receives came
from two distinct sends --- is proved for a FIFO link and rejected for a lossy
one:

\begin{center}\small
\begin{tabular}{@{}p{0.46\linewidth}p{0.30\linewidth}p{0.16\linewidth}@{}}
\toprule
\textbf{\sffamily\scriptsize\MakeUppercase{Claim}} &
\textbf{\sffamily\scriptsize\MakeUppercase{Where}} &
\textbf{\sffamily\scriptsize\MakeUppercase{Outcome}} \\
\midrule
two receives come from two distinct sends & \code{heartbeat.rs}, per-link FIFO
& verifies \\
\addlinespace
the same claim, verbatim & \code{counterexamples/}\newline\code{lossy\_counting.rs}
& rejected \\
\bottomrule
\end{tabular}
\end{center}

A protocol that needs the count must re-establish it, by putting sequence
numbers in its messages and discarding what it has already seen. The obligation
appears exactly where the argument would otherwise be unsound.

\subsection*{Many producers, one consumer}

Unordered delivery from several senders is of the first kind, and it constrains
the model rather than the delivery discipline. A single send history with
several concurrent appenders cannot be exclusively owned, and without exclusive
ownership \code{send} is not a left mover.

So a bag is modelled as a \emph{family} of channels, one per producer. Each
producer exclusively owns its own history, ownership and left-moverness are
restored, and per-producer order is retained --- which a real multi-producer
channel also retains. ``Any order'' moves to the receive side, where it belongs:

\srcname{src/tok.rs}
\begin{lstlisting}
pub fn recv_any<M, Inv: NetInv<M>>(
    rxs: &Vec<Receiver<M>>, ...,
    Tracked(map): Tracked<&mut MapToken<ChanId, Seq<M>, NetSM::recvd<M, Inv>>>,
) -> (res: (usize, M, Tracked<NetSM::was_sent<M, Inv>>))
\end{lstlisting}

It blocks on a family of channels and reports which one produced a message. The
consumer in \code{src/examples/collector.rs} takes whatever arrives first, from
whichever producer, and learns from the yield invariant that it is well formed.

% ===========================================================================
\section{Several threads in one process}\label{sec:threads}

A thread in this development is not a process; it is a unit of sequential
control that owns endpoints. A coordinator that spawns one thread per
participant is ordinary, and needs nothing new.

Sibling threads of one process are, as far as the proof is concerned,
unrelated threads: each preserves the yield invariant, each tolerates the
others' interference, and none mentions the others. What has to be arranged is
not sharing a token but \emph{dividing} one, since ownership is what makes the
framing facts of Section~\ref{sec:net} available at all.

\srcname{src/examples/multithread.rs --- forking and joining}
\begin{lstlisting}
let tracked s0 = sent_map.remove(rsp_chan(0));
let tracked r0 = recvd_map.remove(req_chan(0));

let h0 = vstd::thread::spawn(move || -> (res: (Tracked<...>, Tracked<...>))
    ensures res.0@.key() == rsp_chan(0), ...
    { ... participant_fanout(0, &rx0, &tx0, vote0, ...); (Tracked(ss), Tracked(rr)) });

// ... join, then put the endpoints back:
proof { sent_map.insert(a.0.get()); recvd_map.insert(a.1.get()); }
\end{lstlisting}

Removing a token from a map and inserting it again are proved operations of the
Verus token library, and disjointness holds by construction: what has been
removed is no longer there. Thread creation needs nothing special either, since
a token is an ordinary \code{tracked} value and moves into a spawned closure
like any other value.

Two limitations are worth naming. Sibling threads cannot share a receive
endpoint, which is the single-consumer discipline applied consistently. And
threads of one process sharing mutable state \emph{outside} the network --- a
counter behind a mutex --- is not modelled here at all; that is ordinary
shared-memory concurrency and is orthogonal to the channel reasoning.

% ===========================================================================
\section{Layers}\label{sec:layers}

A protocol proved correct against its own state machine is still stated in terms
of channels and message histories. A specification is usually not: two-phase
commit, seen from above, is a record of who has voted, and a leader election is
an outcome. Layered refinement connects the two.

A \emph{layer} is a gated transition system, and nothing more:

\srcname{src/layer.rs}
\begin{lstlisting}
pub trait Spec {
    type S;   type A;   type M;
    spec fn inv(s: Self::S) -> bool;
    spec fn gate(a: Self::A, req: Self::M, s: Self::S) -> bool;
    spec fn step(a: Self::A, req: Self::M, resp: Self::M, s0: Self::S, s1: Self::S) -> bool;
}
\end{lstlisting}

It says nothing about networks, so an abstract layer may be stated over whatever
state is natural for it. A refinement relates two layers:

\srcname{src/layer.rs}
\begin{lstlisting}
pub trait Refines<Lo: Spec, Hi: Spec> {
    spec fn abs_state(s: Lo::S) -> Hi::S;
    spec fn abs_act(a: Lo::A) -> Option<Hi::A>;    // None: refines skip
    spec fn abs_msg(m: Lo::M) -> Hi::M;

    proof fn lemma_gate(a, req, s) requires Lo::inv(s), Hi::gate(...) ensures Lo::gate(...);
    proof fn lemma_step(a, req, resp, s0, s1) requires Lo::inv(s0), Lo::step(...) ensures ...;
}
\end{lstlisting}

\code{abs\_state} is an Abadi--Lamport refinement mapping: the two layers may
differ in state, actions and messages. \code{abs\_act} returning \code{None}
means the action is invisible from above, which most internal steps of a real
implementation are. The gate condition is Civl's, and runs in the direction one
might not expect: the \emph{abstract} action must fail more often, so that
wherever the abstraction is safe the implementation is.

Both refinement lemmas may assume the concrete invariant, because refinement is
only ever appealed to at reachable states. This matters in practice. In
\code{src/examples/layers\_demo.rs} the abstraction retains only the last sequence
number while the concrete gate speaks of every element of the history; bridging
the two requires knowing the history is increasing, which is exactly the
invariant.

\subsection*{The bottom of the stack is derived, not restated}

A stack has to start somewhere, and the tempting thing is to write the bottom
layer by hand: restate the protocol's gate as \code{gate} and its transition as
\code{step}. That proves refinement about a \emph{copy} of the protocol, with
nothing checking that the copy agrees with the protocol the tokens verify.

\srcname{src/layer.rs}
\begin{lstlisting}
pub trait Layered<M> : NetInv<M> {
    type St;
    spec fn st_inv(s: Self::St) -> bool;
    spec fn st_send(s0: Self::St, s1: Self::St, c: ChanId, q: Seq<M>, m: M) -> bool;
    spec fn st_recv(...) -> bool;
    proof fn lemma_send_gate_st(...);   // the state gate is the machine's own require
}

impl<M, T: Layered<M>> Spec for BottomLayer<(M, T)> { ... }
\end{lstlisting}

A protocol names its machine's own state, invariant and transition relations ---
one line each, pointing at what the state machine library already generated ---
and \code{BottomLayer<T>} is those. There is no second definition to keep in
step. Changing the protocol's gate changes the refinement obligation, which is
the property a hand-written bottom layer does not have.

\subsection*{An expressiveness note}

\code{abs\_act} is a function of the action alone, so an action cannot be visible
in some states and invisible in others. State-dependent invisibility belongs on
the abstract side instead, where the abstract transition simply admits the
identity: \code{HbTop::step} is \emph{the count grows, or nothing changed}.
\code{None} is for actions that are always internal.

% ===========================================================================
\section{What is assumed}\label{sec:assumed}

Verus checks that the obligations above are discharged. It does not, and cannot,
check the reduction theory that makes those obligations sufficient: there is no
formal execution semantics inside the Rust program to state Lipton's theorem
against. So the boundary has to be drawn explicitly.

\begin{center}
\small
\begin{tabular}{@{}p{0.34\linewidth}p{0.60\linewidth}@{}}
\toprule
\textbf{\sffamily\scriptsize\MakeUppercase{Assumed}} &
\textbf{\sffamily\scriptsize\MakeUppercase{Why}} \\
\midrule
\code{send\_general}, \code{recv}, \code{recv\_any}
& The channel primitives. Their specifications are the state machine's
  transitions; the runtime that implements them is not verified.
  \code{send\_general} carries the witnesses for whatever messages the one
  being sent must point at; \code{send} and \code{send\_caused} are verified
  wrappers over it, not further assumptions. \\
\addlinespace
\code{make\_endpoints}
& Fabricates two endpoint handles for a channel identifier. Touches no ghost
  state and is protocol-independent, so creating channels costs no per-protocol
  assumption. \\
\addlinespace
\code{abort\_on\_panicked\_child}
& Diverges, so any postcondition holds. Used when a spawned thread panics and
  its tokens are therefore gone. \\
\addlinespace
Per-protocol configuration
& Four small facts, one per protocol: how many participants, nodes or producers
  there are, and (for the ring) that node identifiers are distinct. \\
\addlinespace
Lipton's theorem
& That pairwise commutativity licenses treating a block as atomic. Its side
  conditions are discharged (Section~\ref{sec:movers}); the theorem is not. \\
\bottomrule
\end{tabular}
\end{center}

Nine \code{external\_body} declarations in total. Note what is \emph{not} on the
list, and was on the corresponding list of an earlier version of this
development: channel freshness, channel distinctness, the reduction rule for
executing a handler at the call site, dividing and recombining ownership across
threads, and the abstracted receive. Each of those is now proved.

The per-protocol residue is worth looking at, because it is small and it is not
about channels:

\begin{lstlisting}
twophase, twophase_fanout:   n_parts() >= 0
collector:                   n_producers() > 0
chang_roberts:               n_nodes() >= 1, and node identifiers are distinct
\end{lstlisting}

These are facts about a deployment. Everything a protocol used to assume about
its own channel naming --- that participants have distinct channels, that a
request channel is never a reply channel, that a minted channel is fresh --- now
follows from how names are built.

\subsection*{Where the assumptions could be reduced further}

The remaining reduction assumption is the substantial one. Building a shallow
trace semantics as specification-level data and proving the reduction theorem
about it would relocate the trust from \emph{reduction is sound}, a claim about
all executions, to \emph{these three functions implement this transition
relation}, which is smaller and reviewable. That is the natural next step, along
with the inductive rule needed for the quorum-gathering protocols of
Section~\ref{sec:fanout}.

One obligation is currently discharged by review rather than by the verifier: a
function body is not tied to a declared action of the protocol. In the earlier
design this connection was written out by hand for one protocol; carrying it
over needs the global state that the token representation deliberately hides,
and is unfinished.

% ===========================================================================
\section{Summary}\label{sec:summary}

For a reader who wants the short version of what writing a protocol involves:

\begin{enumerate}[leftmargin=1.4em,itemsep=3pt,topsep=3pt]
\item Declare the message type and the protocol's parameters.
\item Implement \code{NetInv}: the gate, the guarantee about a message that was
  sent, the delivery discipline, and --- only for a guarantee no single message
  can express --- the \code{extra} invariant and its three lemmas. A guarantee
  that relates messages on \emph{different} channels additionally needs
  \code{needs\_cause}, \code{caused\_by}, \code{cause\_gives} and one lemma
  (Section~\ref{sec:provenance}); all four default to saying nothing.
\item Write the protocol as ordinary sequential Rust, using the library's
  \code{send} and \code{recv} --- or \code{send\_caused}, where a message must
  point at another, or \code{send\_general}, where it must point at several ---
  and passing the state machine instance and the relevant channel token.
\item Implement \code{DetDelivery} if the protocol makes remote calls, and
  \code{Layered} if it is to sit at the bottom of a refinement stack.
\end{enumerate}

There is no mover annotation, no footprint, no commutativity obligation and no
account of what other threads may do. Those obligations have no counterpart
here: the facts they used to establish are consequences of who owns what, and
the mistakes they used to catch cannot be written down.

The cost, stated fairly, is notation. Every function that touches the network
carries the state machine instance and a channel token, together with
preconditions relating them. Code that holds many channels moves tokens in and
out of a map around each operation.

The limit of the approach is that the network is exactly \code{sent},
\code{recvd}, the record of what was sent, and the allocator. A protocol needing
ghost state of its own declares a separate state machine, which composes with
this one at no cost because the two are independent instances. What is not
expressible is protocol state that must be \emph{related} to network state; that
would have to live in one machine. Relating messages on different channels to
each other, which was outside the model in an earlier version of this
development, is now within it, and at no cost in assumptions: the send
primitive was generalised rather than duplicated.

% ===========================================================================
\appendix
\section{Verus notation}\label{sec:verus}

The development is written in Rust; Verus adds specifications and proofs to Rust
and discharges them with an SMT solver. This appendix collects what a reader who
knows Rust but not Verus needs in order to read the code in the body.

\begin{tcolorbox}[colback=sunken,colframe=rulegrey,boxrule=0.4pt,arc=1pt,
                  left=8pt,right=8pt,top=6pt,bottom=6pt,breakable]
{\sffamily\scriptsize\color{neutralmv}\MakeUppercase{Verus notation}}
\par\vspace{5pt}
\begin{description}[leftmargin=!,labelwidth=6.2em,font=\ttfamily\small,
                    itemsep=2pt,parsep=0pt,topsep=2pt]
\item[spec fn] A mathematical function, usable only in specifications. Erased
  before compilation.
\item[proof fn] A lemma. Its \code{requires}/\code{ensures} are its statement;
  its body is its proof. Erased before compilation.
\item[open spec fn] A specification function that can be called in proofs, but not in executable code.
The difference between ``open spec fn'' and ``spec fn'' is that the former can be called in 
proofs, while the latter can only be called in specifications.
\item[fn] Executable code. This is what actually runs.
\item[requires\slash ensures] Pre- and postconditions.
\item[old(x), final(x)] For a parameter passed by mutable reference, its value
  before and after the call. Both appear in postconditions.
\item[Ghost<T>] A value of type \code{T} that exists only in the proof and
  occupies no space at runtime. For a value \code{x} of type \code{Ghost<T>},
we write \code{x@} to read the underlying value.
\item[Tracked<T>] Like \code{Ghost}, but subject to Rust's ownership rules: it
  cannot be duplicated or discarded silently.
\item[Seq, Set, Map] Mathematical sequences, sets and maps. Unbounded;
  specification-only.
\item[\&\&\&, |||] Bulleted conjunction and disjunction, for readable
  multi-line formulas.
\item[x->Variant\_0] Projects the first field of an enum variant. Total:
  applying it to the wrong variant is well-typed and yields an unspecified
  value.
\item[uninterp] An uninterpreted specification function: a parameter of the
  protocol, constrained only by whatever axioms are stated about it.
\item[external\_body] Marks a function whose body is \emph{not} verified. Its
  specification is assumed. Every one of these is an axiom.
\item[tracked x] A proof-mode value subject to ownership: it can be moved and
  consumed, but not copied.
\end{description}
\end{tcolorbox}

One Rust feature matters to the proof rather than to the syntax. Values have
unique owners, and a value that is not \code{Clone} cannot be duplicated. The
development gives \code{Receiver} no \code{Clone} implementation, so a receive
endpoint is owned by exactly one thread, and Rust's type system discharges the
single-consumer condition the proof needs at no cost. Section~\ref{sec:net}
applies the same idea to the network state itself.

\subsection*{Tokenized state machines}

The one Verus library the development relies on is
\code{tokenized\_state\_machine!}. Section~\ref{sec:comp} says why this shape of
declaration is useful; this appendix says what it looks like and what it
generates.

The macro takes a transition system whose state is divided into named
\emph{fields}, each annotated with a \emph{sharding strategy} saying how that
field is divided among its owners. Here is a complete example: a counter that
several threads may increment, where each holds a share of the total.

\srcname{a small tokenized state machine}
\begin{lstlisting}
tokenized_state_machine!{
    Counter {
        fields {
            #[sharding(variable)] pub total: nat,   // one exclusive token
            #[sharding(count)]    pub units: nat,   // splittable into shares
        }

        #[invariant]
        pub spec fn agree(&self) -> bool { self.total == self.units }

        init!{
            boot() { init total = 0; init units = 0; }
        }

        transition!{
            incr() {
                update total = pre.total + 1;       // change an exclusive field
                add    units += (1);                // mint one share
            }
        }

        #[inductive(boot)]  fn boot_inductive(post: Self) { }
        #[inductive(incr)]  fn incr_inductive(pre: Self, post: Self) { }
    }
}
\end{lstlisting}

From this the macro generates:

\begin{itemize}[leftmargin=1.2em,itemsep=3pt,topsep=3pt]
\item a type \code{Counter::Instance}, identifying one running copy of the
  machine, together with \code{Instance::boot()}, which returns the instance and
  \emph{every token of the initial state};
\item a token type per field --- \code{Counter::total}, \code{Counter::units}
  --- whose values are \code{tracked}, so they obey Rust's ownership rules;
\item a proof method per transition, here \code{instance.incr(...)}, which
  consumes the tokens the transition removes and returns those it adds;
\item a proof obligation per \code{\#[inductive(..)]} stub, discharged by the
  bodies above, that the transition preserves every \code{\#[invariant]}.
\end{itemize}

The strategies used in this development are \code{variable} (a single exclusive
token for the whole field), \code{map} (one exclusive token per key),
\code{constant} (a value fixed at initialisation, readable from the instance),
and \code{persistent\_set} (duplicable tokens that are never retracted). Within
a transition body:

\begin{center}\small
\begin{tabular}{@{}p{0.26\linewidth}p{0.68\linewidth}@{}}
\toprule
\textbf{\sffamily\scriptsize\MakeUppercase{Clause}} &
\textbf{\sffamily\scriptsize\MakeUppercase{Meaning}} \\
\midrule
\code{require(P)} & The transition may only be taken where \code{P} holds.
  Whoever performs it must establish \code{P}. \\
\addlinespace
\code{remove f -= [k => v]} & Consume the token for key \code{k}, whose value is
  \code{v}. The performer must hold it. \\
\addlinespace
\code{add f += [k => v]} & Produce a token for \code{k}. For a map, this
  carries an obligation that \code{k} was absent. \\
\addlinespace
\code{have f >= \{x\}} & Read a persistent token without consuming it. \\
\addlinespace
\code{update f = e} & Change a \code{variable} field. \\
\addlinespace
\code{assert(P)} & A proof obligation discharged using the invariants. \\
\addlinespace
\code{birds\_eye let x = e} & Bind a value read from the whole state, including
  fields the performer does not own. The binding is not available to the caller,
  so only what is proved \emph{without} mentioning it is returned
  (Section~\ref{sec:provenance}). \\
\bottomrule
\end{tabular}
\end{center}

A \code{property!} block has the same shape as a transition but changes nothing.
It is how a thread converts tokens it holds into a consequence of the machine's
invariants --- the mechanism Section~\ref{sec:yield} uses to read a yield
invariant off a witness.

Two consequences are worth keeping in mind while reading the body. A thread can
affect the state only by performing a declared transition, and only if it holds
the tokens that transition consumes; and the invariants are established once,
over the declared transitions, rather than re-proved at each use.

\vspace{1.5em}
\noindent\textcolor{rulegrey}{\rule{\linewidth}{0.4pt}}
\vspace{0.4em}

{\small\color{comment}
\noindent\textbf{Sources.} The proof rules are Civl's. See Kragl and Qadeer,
\emph{The Civl Verifier} (FMCAD 2021); Kragl, Qadeer and Henzinger,
\emph{Refinement for Structured Concurrent Programs} (CAV 2020); Kragl, Enea,
Henzinger, Mutluergil and Qadeer, \emph{Inductive Sequentialization of
Asynchronous Programs} (PLDI 2020); and Gangamreddypalli, Enea and Qadeer,
\emph{Reduction for Structured Concurrent Programs} (ESOP 2026).

\vspace{0.4em}
\noindent The library is \code{src/tok.rs} and \code{src/layer.rs}. The
protocols are in \code{src/examples/}; \code{counterexamples/} holds protocols
that must fail to verify. Checked with Verus 0.2026.08.15.
}

\end{document}
